\documentclass{article}
\usepackage{debug/code}

\usepackage{debug/de}
\usepackage{debug/math}
\usepackage{debug/boxes}
\usepackage{debug/diagrams}
\usepackage{debug/anki}
\usepackage{debug/laws}
%\usepackage{release/de}
%\usepackage{release/math}
%\usepackage{release/boxes}
%\usepackage{release/diagrams}
%\usepackage{release/anki}
%\usepackage{release/laws}
\begin{document}
\section{Globals}
\subsection{Deutsch}
\zb, \etc, \idr, \textquote{text}

\subsection{Gesetze}
\art{1} \para{1} \absa{1} \satz{1} \nr{1} \lit{a} \gesetz{BGB}

\GG

\subsection{Mathe}
\subsubsection{Konstanten}
\e

\subsubsection{Mengen}
Die Zahlenmengen \tN, \tNz, \tZ, \tQ, \tI, \tR und \tC. Die Lösungsmenge \tL.
\[
 0 \in \N, \Nz, \Z, \Q, \I, \R, \C
\]

\subsection{Geometrie}
\[
 90\dgr
\]

\subsubsection{Formattierung}
\[
 1 \stext{und} 2, \dtext{und} \e \dvert 3 \ddtext{und} 4 * 5
\]
\begin{mathdescription}
 \mathitem{Text 1} 1 \\
 \mathitem{Text 2} 1
\end{mathdescription}

\section{Code}
\begin{sourcecode}{python}
def fib(n):
    if n <= 1:
        return n
    return fib(n - 1) + fib(n - 2)
\end{sourcecode}
\begin{slsourcecode}{java}
public static void main(String[] args)
\end{slsourcecode}

\subsection{Boxes}
\ztodo{Hello World}

\subsection{Diagramme}
\subsubsection{Farben}
\tikz{
 \draw[fg] (0, 0) -- (1, 1);
 \draw[bg] (0, 0.1) -- (1, 1.1);
 \draw[hl] (0, 0.2) -- (1, 1.2);
 \draw[hl2] (0, 0.3) -- (1, 1.3);
}
\subsubsection{Koordinaten System}
\tikz{
 \coordsys
}
\subsubsection{Koordinaten System (3)}
\tikz{
 \coordsysc{3}
}

\section{Anki}
\ankishow
\ankiprefix{prefix}
\ankitopic{topic}
\ankinote{id}{front}{back}
\end{document}